% ch-a6-02.tex —— 第 6 篇（Feng–Harris）续：PDF 496–507 / 印刷 483–494

\ysrc{483}\begin{thm}{2.6.2}（Avramov，[Avr99]）设 $f\colon A\to B$ 是 Noether 环的同态。则 $f$ 是 LCI 当且仅当 $\mathbf L_f$ 能由一个上同调集中在次数 $[-1,0]$ 中的复形表示。
\end{thm}

\begin{defn}{2.6.3}动态交换环（animated commutative ring）的态射 $f\colon\mathcal A\to\mathcal B$ 称为拟光滑的（quasismooth），若 $\mathbf L_f$ 的 Tor-幅度落在 $[-1,0]$ 中，即 Zariski 局部地在 $\mathcal B$ 上能由位于次数 $-1,0$ 的向量丛复形表示。（这里我们采用上同调编号。）
\end{defn}

更一般地，导出叠的态射 $f\colon X\to Y$ 拟光滑，若 $\mathbf L_f$ 存在，且其 Tor-幅度落在 $[-1,\infty)$ 中，即局部地能由某个 $n<\infty$ 的、位于次数 $[-1,n]$ 的完美复形表示。

若 $f$ 拟光滑，则称 $\chi(\mathbf L_f)$ 为 $f$ 的相对维数。

定理 2.6.2 说明了我们为何说拟光滑性是「LCI」的导出推广。这一性质使一些重要构造成为可能，例如相交理论中「Gysin 拉回」的构造。此外，在某种意义上，自然界中遇到的拟光滑态射比 LCI 态射更常见；这实际上是导出代数几何进入计数代数几何的主要动机之一。例如，知道定义之后便容易证明：

\begin{prop}{2.6.4}拟光滑态射对复合与导出基变换封闭。
\end{prop}

作为对照，注意 LCI 态射不被基变换保持。

\begin{exam}{2.6.5}任何有限型仿射概形都能实现为某个拟光滑导出仿射概形的经典截断，理由很平常。事实上，任取一个展示 $R=\Z[x_1,\dots,x_n]/(f_1,\dots,f_m)$，下图导出纤维积的经典截断同构于 $\Spec R$：
\[
\begin{tikzcd}[column sep=large, row sep=large]
& \Spec \Z[x_1,\dots,x_n] \arrow[d, "{(f_1,\dots,f_m)}"] & \\
\Spec \Z \arrow[r, "0"] & \Spec \Z[y_1,\dots,y_m] &
\end{tikzcd}
\]
这说明拟光滑这一性质对底层的经典截断不施加任何限制。
\end{exam}

若 $f\colon X\to Y$ 有切复形 $\mathbf T_f$，则可用它重新表述拟光滑性：$f$ 拟光滑当且仅当 $\mathbf T_f$ 能由 Tor-幅度落在 $(-\infty,1]$ 中的完美复形表示。

\begin{exam}{2.6.6}由例 2.5.1，Zariski 局部同构于光滑概形的导出纤维积的导出概形是拟光滑的。反之亦然。
\end{exam}

\section{与朗兰兹纲领相关的导出模空间}

导出代数几何在朗兰兹纲领中的相关性通过导出模空间体现。本节概述一些例子，讨论如何看待它们，并提示它们为什么有用。

\ysrc{484}\subsection{隐藏光滑性哲学}

从历史上看，导出模空间这一概念在 20 世纪 80 年代就已被预见——甚至早于导出代数几何的出现——按我们的理解，预见者是 Beilinson、Deligne、Drinfeld 与 Kontsevich。他们强调了一条原则，称为隐藏光滑性哲学（hidden smoothness philosophy），它预测自然模空间应当具有拟光滑（并且具有「预期维数」）的导出版本。

这条原则的背景是这样的：在经典代数几何的某些领域中，人们经常遇到既不是 LCI、维数也不「对」的经典模空间，也就是经典概形或叠。在这种情形下，隐藏光滑性哲学预测，它们是某些具有「正确」几何性质的导出模空间的经典截断。于是，找出这些导出强化的定义便可能有用。

本节从这一角度讨论导出模空间：在试图把一个经典模空间「升级」为导出模空间时，应当寻找什么。随后的几节将（但愿能）说明这么做为什么可能有用。

据我们所知，隐藏光滑性哲学的动机是经验性的，我们将通过例子来说明它。

\begin{exam}{3.1.1}（Betti 局部系统的模空间）我们先非正式地讨论导出代数几何中一个稍早的例子，它与（Betti）几何朗兰兹 [BZN18] 相关：光滑射影（连通）曲线 $C/\C$（亏格 $g$）上（Betti）局部系统的模空间。选定基点之后，可以把 Betti 局部系统视为 $C$ 的基本群的表示。不选基点也有更典范的表述方式，但为了具体起见，我们选一个基点 $c_0\in C$，于是 $\pi_1(C,c_0)$ 有标准展示
\begin{equation}
\pi_1(C,c_0)\approx\langle a_1,\dots,a_g,b_1,\dots,b_g:\prod_{i=1}^{g}[a_i,b_i]=1\rangle.
\tag{3.1.1}\label{eq:a6-02-3-1-1}
\end{equation}
这个展示给出了写下 $C$ 上秩 $n$（Betti）局部系统模空间的一个合理猜测，即取纤维积
\begin{equation}
\begin{tikzcd}[column sep=large, row sep=3.2em]
& & \mathrm{GL}_n^{2g} \arrow[d, "{\displaystyle\prod_{i=1}^{g}[x_i,y_i]}"] \\
e \arrow[rr] & & \mathrm{GL}_n
\end{tikzcd}
\tag{3.1.2}\label{eq:a6-02-3-1-2}
\end{equation}
再对 $\mathrm{GL}_n$ 的共轭作用取商。由于图 (3.1.2) 中所有空间都光滑（特别地，底部的水平箭头是正则嵌入），这提示局部系统的模空间可能是局部完全交，其维数等于 $2g\dim\mathrm{GL}_n-2\dim\mathrm{GL}_n$。然而事实并非总是如此，因为切出 $\{e\}\hookrightarrow\mathrm{GL}_n$ 的正则序列未必能拉回到 $\mathrm{GL}_n^{2g}$ 中的正则序列。

不过，如果不在经典概形中取图 (3.1.2) 的纤维积，而是取导出纤维积，则由命题 2.6.4，所得的导出概形拟光滑，并且具有「预期」维数 $2(g-1)\dim\mathrm{GL}_n$。
\end{exam}

许多例子中，隐藏光滑性哲学之所以成立，是因为对自然出现的模空间可以找到自然的（局部）展示，

\ysrc{485}与 (3.1.2) 类似。更准确地说，拟光滑性与维数条件由形变理论提示，这将在下文说明。

本节解释求解下述问题的启发式方法：
\begin{quote}
给定一个经典模空间 $M$，找出以它为经典截断的「正确」导出模空间 $\mathcal M$。
\end{quote}
按定义，经典模空间是定义在交换环范畴上的一类函子，而导出模空间是定义在单纯交换环范畴上的一类函子，因此这个问题的内容是：以「好的」方式延拓该函子的定义域。（目标范畴也应延拓，从集合或群胚延拓到单纯集。）类比导出函子的概念，我们非正式地把这个问题称为「导出经典模空间 $M$」。

当然，这个问题极不确切：一个给定的经典模空间可以实现为许多不同的（拟光滑）导出模空间的经典截断（见例 2.6.5），正如一个给定的既约概形可以实现为许多不同概形的底层既约概形。所以，我们所说的「正确」的导出模空间是启发式的，由具体问题的需要决定。实践中，可以通过希望得到拟光滑的 $\mathcal M$、甚至希望得到切复形特别的 $\mathcal M$ 来使问题变得刚性。

本节描述这一过程在实践中通常展开的几种模式与启发式，并用两个主要例子说明：Galatius–Venkatesh 的导出伽罗瓦形变环 [GV18]，以及 Feng–Yun–Zhang 的导出 Hitchin 叠 [FYZ21]。

\subsection{导出模空间的启发}

设我们从一个经典模空间 $M$ 出发，想把它提升为导出模空间 $\mathcal M$。下面的纲要描述了寻找「正确」导出强化 $\mathcal M$ 的最标准模式。

设我们处于下述情形：
\begin{itemize}
\item 可以算出 $M$ 在点 $m\in M$ 处的切空间，并且它有某种自然上同调群的解释。此外，自同构落在「前一个」上同调群中，障碍落在「后一个」上同调群中。
\end{itemize}
那么应当尝试构造 $\mathcal M$，使它的切复形（参见 §2.4.3）是计算前面列出的那些上同调群的上同调复形。

我们将给出一些具体例子来说明这一点。不过在此之前，先讨论这一过程何时不必要。

\subsubsection{经典性的启发}

本小节给出一些指导原则，用于判断何时应当期待一个不能由经典概形或叠表示的导出强化。基本哲学由下面这句口号概括。
\begin{quote}
\textbf{Slogan.} 若一个模空间是「正确」维数的 LCI，则它不需要导出。
\end{quote}
「正确」维数这一概念本身是启发式的，它可能来自 (3.1.2) 那样的相交理论性设置，也可能来自计算预期余切复形的欧拉示性数。

\ysrc{486}更准确地说，一个模空间「不需要导出」是指：「正确」的导出模空间应当同构于它的经典截断。

这条原则的依据如下。设 $\mathcal M$ 是导出概形，经典截断为 $M\hookrightarrow\mathcal M$。按隐藏光滑性哲学，设 $\mathcal M$ 是「正确」维数的拟光滑空间，而该维数在 $m\in\mathcal M$ 处就是余切复形在该点的欧拉示性数 $\chi(\mathbf L_{\mathcal M,m})$。

\begin{lem}{3.2.1}典范映射 $\iota\colon M\to\mathcal M$ 诱导同构 $H^0(\iota^*\mathbf L_{\mathcal M})\xrightarrow{\sim}H^0(\mathbf L_{\mathcal M})$ 与满射 $H^{-1}(\iota^*\mathbf L_{\mathcal M})\twoheadrightarrow H^{-1}(\mathbf L_{\mathcal M})$。
\end{lem}

我们不给引理 3.2.1 的证明，它可在 [GV18, §7.2] 中找到；但至少给出一些直观。在单纯交换环的层面上，经典截断是通过在次数 $\ge2$ 上添加单形以「杀掉」高阶同伦群来实现的。在拓扑学中，Hurewicz 定理粗略地说：拓扑空间的映射诱导低次同伦群的同构，当且仅当它诱导低次同调群的同构。引理 3.2.1 来自单纯交换环的类似估计。

\begin{rem}{3.2.2}当 $\mathcal M$ 是导出叠时，引理 3.2.1 有变体：映射 $\iota^*\mathbf L_{\mathcal M}\to\mathbf L_{\mathcal M}$ 在次数 $\ge0$ 上诱导同构，在次数 $-1$ 上诱导满射。
\end{rem}

现在设 $M$ 是「正确」维数的 LCI，这里把「正确维数」理解为 $\chi(\mathbf L_{\mathcal M})$。则由定理 2.6.2，$M$ 的余切复形的 Tor-幅度落在 $[-1,0]$ 中，而其欧拉示性数与 $\mathbf L_{\mathcal M}$ 的相同。于是引理 3.2.1 迫使 $\iota^*\mathbf L_{\mathcal M}\to\mathbf L_{\mathcal M}$ 是同构，再由引理 2.4.5 可知映射 $M\to\mathcal M$ 本身是同构。

我们没有把上述论证变成一般的正式陈述，因为「正确」维数是由启发式确定的：证明中用到的维数是 $\chi(\mathbf L_{\mathcal M})$，但这是循环的，因为它预设了 $\mathcal M$ 的存在。实践中，「正确」与否由应用是否成功来判定。我们要提醒：自然界中确实会出现 $M$ 是 LCI（甚至光滑）但维数并不「正确」的例子，例如 [GV18] 中研究的（带晶体条件的）导出伽罗瓦形变环。

\begin{rem}{3.2.3}实践中，要验证模空间 $M$ 是正确维数的 LCI 可能很难，除非 $M$ 还是正确维数的光滑空间。这是因为 $M$ 的余切复形通常不易计算，除非障碍群消没（此时它只是一个层，而 $M$ 光滑）。相反，对导出模空间 $\mathcal M$ 而言，其余切复形通常有具体的描述，因为它可以用「高阶导出对偶数」来解释；关于这一点，§9.1 有更多讨论。
\end{rem}

\subsubsection{经典模空间的贯穿例子}

在 §3.2 开头的纲要中我们说：
\begin{quote}
我们从一个经典模空间 $M$ 出发，想把它提升为导出模空间 $\mathcal M$。
\end{quote}
下面列出一些例子，它们将作为反复回到的贯穿例证。

\ysrc{487}\begin{exam}{3.2.4}（Betti 局部系统的模空间）设 $C$ 是亏格 $g$ 的光滑射影曲线，$G$ 是 $\C$ 上的约化群。受例 3.1.1 启发，我们定义 $C$ 上 Betti $G$-局部系统的模叠为如下得到的叠：取纤维积
\begin{equation}
\begin{tikzcd}[column sep=large, row sep=3.2em]
\mathrm{M}^{\mathrm{Betti},\square}_G \arrow[r] \arrow[d] & G^{2g} \arrow[d, "{\displaystyle\prod_{i=1}^{g}[x_i,y_i]}"] \\
e \arrow[r] & G
\end{tikzcd}
\tag{3.2.1}\label{eq:a6-02-3-2-1}
\end{equation}
所得的 $\mathrm{M}^{\mathrm{Betti},\square}_G$，再定义 $\mathrm{M}^{\mathrm{Betti}}_G:=\mathrm{M}^{\mathrm{Betti},\square}_G/G$，即对共轭 $G$-作用取商。

这个模空间是 Betti 几何朗兰兹猜想的「谱侧」的研究对象，该猜想由 Ben-Zvi–Nadler 提出 [BZN18]。
\end{exam}

\begin{exam}{3.2.5}（整体无限制伽罗瓦形变环）设 $F$ 是整体域，$\mathcal O_F\subset F$ 是其整数环。设 $S$ 是 $\mathcal O_F$ 的有限素位集，$\mathcal O_F[1/S]$ 是 $\mathcal O_F$ 在 $S$ 处的局部化。设 $G$ 是分裂约化群，$\overline\rho\colon\pi_1(\mathcal O_F[1/S])\to G(\F_p)$ 是伽罗瓦表示。\footnote{我们略去了基点的选取；更好的视角见 §5.2.3。}我们定义 $\overline\rho$ 的（整体无限制的）伽罗瓦形变函子为 $\Z_p$-代数范畴上的叠 $\mathrm{M}^{\mathrm{Gal}}_G$：其对象是带增广 $A\xrightarrow{\epsilon}\F_p$ 的 Artin $\Z_p$-代数 $A$，$A$-点是提升 $\rho\colon\pi_1(\mathcal O_F[1/S])\to G(A)$ 的群胚，这些提升经增广 $\epsilon$ 送到 $\overline\rho$，如下图所示：
\begin{equation}
\begin{tikzcd}[column sep=large, row sep=large]
& G(A) \arrow[d, "\epsilon"] \\
\pi_1(\mathcal O_F[1/S]) \arrow[ur, "\rho"] \arrow[r, "\overline\rho"'] & G(\F_p)
\end{tikzcd}
\tag{3.2.2}\label{eq:a6-02-3-2-2}
\end{equation}
这类形变函子由 Mazur 在 [Maz89] 中引入，是 Taylor–Wiles 开创的自守提升定理的基础；更多内容见 Caraiani–Shin 的文章 [CS25]。（这里「整体」表示我们看的是整体域中某个 $S$-整数环的基本群的表示，而「无限制」指我们没有施加任何局部条件——为了应用到朗兰兹对应，最终我们是希望施加这类条件的。）
\end{exam}

\begin{rem}{3.2.6}$\mathrm{M}^{\mathrm{Gal}}_G$ 是否需要导出，微妙地依赖于 $F$ 与 $G$。我们假设 $\overline\rho$ 绝对不可约且奇（odd），即所有复共轭的迹都取极小值——见 [Gro]。例如，若 $F=\Q$ 且 $G=\mathrm{GL}_1,\mathrm{GL}_2$，则不应期待任何有趣的导出强化；而若 $G=\mathrm{GL}_3,\mathrm{GL}_4,\dots$，则应当期待。若 $F$ 是非全实的数域且 $G$ 半单，则 $\mathrm{M}^{\mathrm{Gal}}_G$ 应当是导出的；反之若 $F$ 是函数域，则由 [dJ01, Gai07]，对任何 $G$，$\mathrm{M}^{\mathrm{Gal}}_G$ 都不必导出。最后我们提一句：若 $G=\mathrm{GL}_1$（且 $p\ne2$），则除非 $F=\Q$ 或 $F$ 是虚二次域，$\mathrm{M}^{\mathrm{Gal}}_G$ 都应当导出；例 5.5.1 对此有更多讨论。
\end{rem}

\begin{exam}{3.2.7}（局部无限制伽罗瓦形变）下文将讨论整体域的（无限制）伽罗瓦形变函子的导出问题。人们也可以问：局部域的（无限制）伽罗瓦形变函子是否也应当导出。\ysrc{488}可以证明它们总是正确维数的 LCI，因此不必导出：$\ell\ne p$ 的情形见 [DHKM20]，$\ell=p$ 的情形见 [BIP23]。相反，带条件的局部形变函子是否应当导出，目前还不清楚，这似乎是导出伽罗瓦形变理论未来的一项重要问题；§9.1 有更多讨论。
\end{exam}

\begin{exam}{3.2.8}（Hitchin 叠）设 $C$ 是 $\F_q$ 上特征 $p>2$ 的光滑射影曲线。对整数 $m,n\ge0$，定义 $\F_q$ 上的 Hitchin 叠 $\mathrm{M}^{\mathrm{Hitch}}_{m,n}$，其 $R$-点是五元组 $(\mathcal E_1,\mathcal E_2,\mathcal F,t_1,t_2)$ 的群胚，其中
\begin{itemize}
\item $\mathcal E_i$ 是 $C_R=C\times_{\F_q}\Spec R$ 上秩 $m$ 的向量丛，
\item $\mathcal F$ 是 $C_R$ 上秩 $n$ 的向量丛，
\item $t_1\in\Hom(\mathcal E_1,\mathcal F)$，$t_2\in\Hom(\mathcal F,\mathcal E_2^\vee)$，其中 $\mathcal E_2^\vee\simeq\mathcal E_2^*\otimes\omega_C$ 是 $\mathcal E_2$ 的 Serre 对偶。
\end{itemize}
这是 [FYZ24, FYZ21] 中的酉 Hitchin 空间 $M$ 在「分裂情形」（酉群退化为一般线性群）下的特化。引入它是为了分析 shtuka 模叠上的特殊闭链，后者出现在算术 theta 级数的函数域版本中。\footnote{空间 $\mathrm{M}^{\mathrm{Hitch}}_{m,n}$ 不是 Hitchin 最初研究的 Higgs 丛空间，而是类似构造的另一些实例，它们在自守表示论中已变得重要——更多讨论见 [Yun18, FW25]。}
\end{exam}

\subsubsection{来自形变理论的提示}

现在设经典模空间 $M$ 已给定。接下来，在 §3.2 中我们说：
\begin{quote}
可以算出 $M$ 在点 $m\in M$ 处的切空间，并且它有某种自然上同调群的解释。此外，自同构落在「前一个」上同调群中，障碍落在「后一个」上同调群中。
\end{quote}
下面解释这一点。关键在于：$M$ 在 $m\in M$ 处的切空间有一阶形变的「具体」描述，这来自 $M$ 作为模空间的定义，以及用对偶数对切空间的解释。

\begin{exam}{3.2.9}（Betti 局部系统的模空间）沿用例 3.2.4 的记号，设 $L$ 是 $C$ 上（Betti）$G$-局部系统。可以把 $L$ 视为 $\mathrm{M}^{\mathrm{Betti}}_G$ 的 $\C$-点。则 $\mathrm{M}^{\mathrm{Betti}}_G$ 在 $L$ 处的切空间是 $H^1(\pi_1(C),\frakg_L)$，其中 $\frakg_L:=L\times^G\frakg$ 是 $C$ 上的局部系统，通过取 $L\times\frakg$ 对对角 $G$-作用的商得到。这个局部系统有时也记作 $\operatorname{Ad}L$ 或 $\frakg_L$。此外，
\begin{itemize}
\item $L$ 的自同构可以典范地等同于 $H^0(\pi_1(C),\frakg_L)$，并且
\item 可以证明一阶形变的障碍落在 $H^2(\pi_1(C),\frakg_L)$ 中。
\end{itemize}
\end{exam}

\begin{exam}{3.2.10}（整体无限制伽罗瓦形变）沿用例 3.2.5 的记号，把 $\overline\rho$ 视为 $\mathrm{M}^{\mathrm{Gal}}_G$ 的 $\F_p$-点。则 $\mathrm{M}^{\mathrm{Gal}}_G$ 在 $\overline\rho$ 处的切空间是 $H^1(\pi_1(\mathcal O_F[1/S]),\frakg_{\overline\rho})$，其中 $\frakg_{\overline\rho}$ 的底向量空间是 $\frakg_{\F_p}$，\ysrc{489}$\pi_1(\mathcal O_F[1/S])$ 经由 $\overline\rho\colon\pi_1(\mathcal O_F[1/S])\to G(\F_p)$ 作用，而 $G(\F_p)$ 经由它在 $\frakg_{\F_p}$ 上的伴随作用。此外，
\begin{itemize}
\item $\overline\rho$ 的自同构可以典范地等同于 $H^0(\pi_1(\mathcal O_F[1/S]),\frakg_{\overline\rho})$，并且
\item 可以证明一阶形变的障碍落在 $H^2(\pi_1(\mathcal O_F[1/S]),\frakg_{\overline\rho})$ 中 [Maz89, Proposition 2]。
\end{itemize}
注意这与例 3.2.9 类似。
\end{exam}

\begin{exam}{3.2.11}（Hitchin 叠）沿用例 3.2.8 的记号，设 $(\mathcal E_1,\mathcal E_2,\mathcal F,t_1,t_2)\in\mathrm{M}^{\mathrm{Hitch}}_{m,n}(k)$，其中 $k/\F_q$ 是域。令 $\mathbf T(\mathcal E_1,\mathcal E_2,\mathcal F,t_1,t_2)$ 为复形
\[
\operatorname{End}(\mathcal E_1)\oplus\operatorname{End}(\mathcal E_2)\oplus\operatorname{End}(\mathcal F)\xrightarrow{d}\Hom(\mathcal E_1,\mathcal F)\oplus\Hom(\mathcal F,\mathcal E_2^\vee),
\]
其中左项位于次数 $-1$，右项位于次数 $0$，微分把 $(A_1,A_2,B)\in\operatorname{End}(\mathcal E_1)\oplus\operatorname{End}(\mathcal E_2)\oplus\operatorname{End}(\mathcal F)$ 送到 $(Bt_1-t_1A_1,\ A_2^\vee t_2-t_2B)\in\Hom(\mathcal E_1,\mathcal F)\oplus\Hom(\mathcal F,\mathcal E_2^\vee)$。于是 $\mathrm{M}^{\mathrm{Hitch}}_{m,n}$ 在 $(\mathcal E_1,\mathcal E_2,\mathcal F,t_1,t_2)$ 处的切空间可典范地等同于 $H^1(C_k,\mathbf T(\mathcal E_1,\mathcal E_2,\mathcal F,t_1,t_2))$。此外，如 [Ngo10, §4.14] 所示，
\begin{itemize}
\item $(\mathcal E_1,\mathcal E_2,\mathcal F,t_1,t_2)$ 的自同构可以典范地等同于 $H^0(C_k,\mathbf T(\mathcal E_1,\mathcal E_2,\mathcal F,t_1,t_2))$，并且
\item 一阶形变的障碍落在 $H^2(C_k,\mathbf T(\mathcal E_1,\mathcal E_2,\mathcal F,t_1,t_2))$ 中。
\end{itemize}
\end{exam}

\subsubsection{导出模叠的构造}

至此我们有经典模空间 $M$，以及对导出模空间 $\mathcal M$ 的（余）切复形的一个候选猜测。最后，在 §3.2 中我们说：
\begin{quote}
然后尝试构造 $\mathcal M$，使其切复形是计算上一步那些上同调群的上同调复形。
\end{quote}
下面用贯穿例子说明这如何展开。

\begin{exam}{3.2.12}（Betti 局部系统的导出模空间）沿用例 3.2.9 的记号。其实这里已经有一处微妙之处：例 3.2.9 中的上同调是以局部系统 $\frakg_L$ 为系数的 $\pi_1(C)$ 的群上同调。它对 $C$ 的依赖只通过基本群 $\pi_1(C)$，而 $\pi_1(C)$ 可以看作只捕捉了 Postnikov 截断 $\tau_{\le1}C$ 的信息；但正确的做法是把 $C$ 的整个「同伦型」纳入进来。这意味着，我们希望局部系统 $L$ 处的切复形是上同调复形 $R\Gamma_{\mathrm{Betti}}(C,\frakg_L[1])$，而不是 $R\Gamma_{\mathrm{group}}(\pi_1(C),\frakg_L[1])$。

我们简要讨论这两个对象的差别。有
\[
H^0(C,\frakg_L)\cong H^0(\pi_1(C),\frakg_L)\qquad\text{以及}\qquad H^1(C,\frakg_L)\cong H^1(\pi_1(C),\frakg_L),
\]
所以差别出现在 $H^2$。进一步，若 $C$ 的亏格 $g\ge1$，则 $C$ 实际上是 $K(\pi_1(C),1)$，于是 $R\Gamma_{\mathrm{Betti}}(C,\frakg_L[1])\cong R\Gamma_{\mathrm{group}}(\pi_1(C),\frakg_L[1])$。因此（在光滑真复曲线中）这一区别只在 $C\cong\PP^1$ 时出现。此时 $\pi_1(C)$ 平凡，故群上同调平凡；但 $\PP^1$ 的高阶 Betti 上同调非平凡（注意 $\pi_1(\PP^1)=0$，故 $L$ 必是「那个」平凡局部系统）。
\end{exam}

\ysrc{490}构造所需导出模叠 $\mathcal M^{\mathrm{Betti}}_G$ 的一种方法是：取 (3.1.2) 的导出纤维积，然后对 $G$ 的（对角）共轭作用取商。为计算所得对象的切复形，可利用切复形的构造尊重同伦极限（例 2.5.1），以及 $C(\C)$ 的同伦型有如下同伦推出的展示：
\begin{equation}
\begin{tikzcd}[column sep=large, row sep=large]
S^1 \arrow[r] \arrow[d] & \bigvee_{i=1}^{2g}(S^1) \arrow[d] \\
\mathrm{pt} \arrow[r] & C(\C)
\end{tikzcd}
\tag{3.2.3}\label{eq:a6-02-3-2-3}
\end{equation}
注意当 $g=0$（即 $C\cong\PP^1$）时，我们对 $\mathcal M^{\mathrm{Betti}}_G$ 的展示是 $G$ 中恒等元的导出自交再对 $G$ 取商，其维数为 $-2\dim G$。$\mathcal M^{\mathrm{Betti}}_G$ 的经典截断就是 $\mathrm{M}^{\mathrm{Betti}}_G\cong[\Spec\C]/G$，这反映 $\PP^1$ 上只有唯一的局部系统（平凡者）。一个更朴素的尝试是把 $\mathrm{M}^{\mathrm{Betti}}_G$ 升级为切复形为 $R\Gamma(\pi_1(C),\frakg_L[1])$ 的导出模空间，那只会得到 $[\Spec\C]/G$：它仍然拟光滑，但维数错误，为 $-\dim G$。

注意，$\mathcal M^{\mathrm{Betti}}_G$ 的切复形可以描述为与 $R\Gamma_{\mathrm{Betti}}(C_R,\frakg_L[1])$ 同构的完美复形，且这一描述对 $L\in\mathcal M^{\mathrm{Betti}}_G(R)$ 具有函子性；这说明 $\mathcal M^{\mathrm{Betti}}_G$ 拟光滑，因为 $C$ 的 Betti 上同调振幅落在 $[0,2]$ 中。

\begin{exam}{3.2.13}（整体无限制伽罗瓦形变的导出）沿用例 3.2.10 的记号。例 3.2.10 的计算提示，我们应当寻找导出叠 $\mathcal M^{\mathrm{Gal}}_G$，使其在 $\rho\in\mathcal M^{\mathrm{Gal}}_G(R)$ 处的切复形自然地（对 $R$）同构于 $R\Gamma(\pi_1(\mathcal O_F[1/S]),\frakg_\rho[1])$。注意，根据例 3.2.12，人们可能期待这个答案应当调整为 $R\Gamma_{\text{ét}}(\mathcal O_F[1/S],\frakg_\rho[1])$。把后者视为「正确」答案或许确实更好，但整体域中的 $S$-整数环是「étale $K(\pi,1)$」，因此两者总是一致的。\footnote{然而，当 $F$ 是 $\F_q$ 上光滑射影曲线 $C$ 的函数域时，研究 $C$ 上局部系统的处处非分歧形变也很自然。此时切复形是 $R\Gamma_{\text{ét}}(C;\frakg_\rho[1])$，它可以有非平凡的 $H^3_{\text{ét}}$。因此在这种情况下，伽罗瓦形变空间不是拟光滑的。}所需的 $\mathcal M^{\mathrm{Gal}}_G$ 就是 Galatius–Venkatesh 的导出伽罗瓦形变函子。我们暂不给出它的构造，那需要一点同伦论的插曲。从精神上说，人们希望把点函子定义为 (3.2.2)，只不过把 $A$ 换成「在 $\F_p$ 上带增广的 Artin 动态交换环」。所得对象应当是函子：从在 $\F_p$ 上带增广的 Artin 动态交换环的范畴，到 anima（亦称「单纯集」「空间」或「$\infty$-群胚」）的范畴。数域上同调的决定属于类域论与 Poitou–Tate 对偶。结论是：若 $p\ne2$，则 $R\Gamma(\pi_1(\mathcal O_F[1/S]),\frakg_\rho)$ 的 Tor-幅度落在 $[0,2]$ 中，因此 $\mathcal M^{\mathrm{Gal}}_G$ 拟光滑。
\end{exam}

\begin{exam}{3.2.14}（导出 Hitchin 叠）沿用例 3.2.11 的记号，我们期待导出 Hitchin 叠 $\mathcal M^{\mathrm{Hitch}}_{m,n}$ 满足：其余切复形拉回到 $(\mathcal E_1,\mathcal E_2,\mathcal F,t_1,t_2)\in\mathcal M^{\mathrm{Hitch}}_{m,n}(R)$ 上自然同构于 $R\Gamma(C_R,\mathbf T(\mathcal E_1,\mathcal E_2,\mathcal F,t_1,t_2))$。这个 $\mathcal M^{\mathrm{Hitch}}_{m,n}$ 是 [FYZ21] 中出现的导出 Hitchin 叠的「分裂」特例。\ysrc{491}它可以作为 [TV08] 意义下的「导出映射叠」导出得到。或者，也可以把它视为与完美复形
\[
\begin{aligned}
(\mathcal E_1,\mathcal E_2,\mathcal F)\in\Bun_{\mathrm{GL}_m}\times\Bun_{\mathrm{GL}_m}\times\Bun_{\mathrm{GL}_n}(R)
\mapsto{} & R\operatorname{Hom}_{C_R}(\mathcal E_1,\mathcal F)\\
&{}\oplus R\operatorname{Hom}_{C_R}(\mathcal F,\mathcal E_2^\vee)
\end{aligned}
\]
相伴的导出向量丛（§2.5.2）。
\end{exam}

\begin{rem}{3.2.15}从这些例子可以看出，隐藏光滑性哲学在数论中出现的空间里之所以成立，理由似乎是数域及相关环的上同调维数低，这使得切复形「很短」。
\end{rem}

\begin{exam}{3.2.16}（谱 Hitchin 叠与谱周期）非拟光滑的导出模空间在朗兰兹纲领中也有意义。[FW25, BZSV24] 研究了一些例子，并论证它们把自守表示论中周期的谱（即伽罗瓦论的）类比几何化了。一个例子是谱 Hitchin 叠 $\mathrm{M}^{\mathrm{spec}}_{m,n}$，它参数化与 $\mathrm{M}_{m,n}$ 的点类似的 $(\mathcal E_1,\mathcal E_2,\mathcal F,t_1,t_2)$，只是把丛换成了局部系统。为了具体，设 $C/\C$ 是光滑射影曲线（不过对有限域上的曲线也可以表述 $\ell$-进变体）。$\mathrm{M}^{\mathrm{spec}}_{m,n}$ 的 $R$-点是三元组\tn{底本此处作「triples（三元组）」，但所列对象共有五个，疑为「tuples（元组）」之误（印刷 491）。}$(E_1,E_2,F,t_1,t_2)$，其中
\begin{itemize}
\item $E_i$ 是 $C$ 上秩 $m$ 局部系统的 $R$-族，
\item $F$ 是 $C$ 上秩 $n$ 局部系统的 $R$-族，并且
\item $t_1\in\Hom(E_1,F)$ 是从 $E_1$ 到 $F$ 的平坦态射，$t_2\in\Hom(F,E_2^*)$ 是从 $F$ 到 $E_2^*$ 的平坦态射。
\end{itemize}
它有一个导出强化 $\mathcal M^{\mathrm{spec}}_{m,n}$，其切复形是以类似例 3.2.11 的复形 $\mathbf T(E_1,E_2,F,t_1,t_2)$ 为系数的 $C$ 的 de Rham 上同调。然而，$C$ 的 de Rham 上同调维数为 $2$，而 $C$ 的凝聚上同调维数为 $1$，故这个切复形在次数 $2$ 上确实有非零上同调，于是 $\mathcal M^{\mathrm{spec}}_{m,n}$ 不拟光滑。事实上，$\mathcal M^{\mathrm{spec}}_{m,n}$ 不拟光滑这一点是 [FW25] 的一个重要方面，该文研究自守周期与谱周期之间的关系。[FW25] 的主要观点之一是：谱 Hitchin 叠的对偶层——它们在某种意义上度量拟光滑性的失效——是自守周期的谱对应物。
\end{exam}

\section{导出特殊闭链与高阶 theta 函数}

我们将说明例 3.2.14 引入的导出 Hitchin 叠在特殊闭链与高阶算术 theta 级数理论中的应用，该理论由 Feng–Yun–Zhang 在一系列文章中发展 [FYZ24, FYZ21, FYZ23]。

这一讨论的背景是 Kudla 纲领，正如 Chao Li 的文章 [Li25] 所讨论的。它涉及所谓的算术 theta 函数，即由特殊闭链拼成的 Fourier 级数。这些函数最初在正交或酉志村簇上研究，但我们将聚焦于函数域情形，即 Feng–Yun–Zhang 的工作。经典背景与动机参见 [Li25]。

\subsection{Hermite shtuka}

我们回忆 [FYZ24, §6,7] 中的一些定义。设 $X$ 是 $\F_q$ 上的光滑射影曲线，$\nu\colon X'\to X$ 是有限 étale 双重覆盖，$X'$ 在 $X$ 上的非平凡自同构记为 $\sigma\colon X'\to X'$。

\ysrc{492}对每个 $r\ge0$ 与 $n>0$，[FYZ24, §6] 与 [FYZ21, §2] 构造了秩 $n$ 酉 shtuka 的模空间，记作 $\mathrm{Sht}^r_{U(n)}$。它们是 Zhiwei Yun 讲义 [Zhi25] 中所讨论空间的变体。存在映射 $\pi\colon\mathrm{Sht}^r_{U(n)}\to(X')^r$，称为「leg 映射」。当 $r=1$ 时，应把 $\pi$ 大致类比为酉志村簇整模型的结构态射。\footnote{据 Yun 解释，后者的更接近的类比是 $r=2$ 时 $\pi$ 在 $X'\times\{\infty\}$ 上的纤维，其中 $\infty\in X'$ 是固定点。}下面概述这些构造。

对测试概形 $S$，关于 $\nu\colon X'\to X$ 的、$X\times S$ 上秩 $n$ 的 Hermite 丛（Hermitian bundle）是指 $X'\times S$ 上秩 $n$ 的向量丛 $\mathcal F$，连同一个同构 $h\colon\mathcal F\xrightarrow{\sim}\sigma^*\mathcal F^\vee$ 使得 $\sigma^*h^\vee=h$。这里 $\mathcal F^\vee=\mathcal F^*\otimes_{\OO_{X'}}\omega_{X'}$ 是 $\mathcal F$ 的 Serre 对偶。我们把 $h$ 称为「Hermite 结构」。

我们用 $\Bun_{U(n)}$ 记 $X$ 上秩 $n$ Hermite 丛的模叠，它把测试概形 $S$ 送到 $X'\times S$ 上带 Hermite 结构的秩 $n$ 向量丛的群胚。

\begin{defn}{4.1.1}设 $r\ge0$ 是整数。Hecke 叠 $\mathrm{Hk}^r_{U(n)}$ 的 $S$-点是下列数据的群胚：
\begin{enumerate}
\item $x'_i\in X'(S)$，$i=1,\dots,r$，其图记为 $\Gamma_{x'_i}\subset X'\times S$。
\item $X'\times S$ 上秩 $n$ 的向量丛序列 $\mathcal F_0,\dots,\mathcal F_r$，各带 Hermite 结构 $h_i\colon\mathcal F_i\xrightarrow{\sim}\sigma^*\mathcal F_i^\vee$。
\item 同构 $f_i\colon\mathcal F_{i-1}|_{X'\times S-\Gamma_{x'_i}-\Gamma_{\sigma(x'_i)}}\xrightarrow{\sim}\mathcal F_i|_{X'\times S-\Gamma_{x'_i}-\Gamma_{\sigma(x'_i)}}$，$1\le i\le r$，与 Hermite 结构相容，且满足下述性质：存在 $X'\times S$ 上秩 $n$ 的向量丛 $\mathcal F^\flat_{i-1/2}$ 与向量丛的图 (4.1.1)，
\begin{equation}
\begin{tikzcd}[column sep=large, row sep=large]
& \mathcal F^{\flat}_{i-1/2} \arrow[dl, "f_i^{\leftarrow}"'] \arrow[dr, "f_i^{\rightarrow}"] & \\
\mathcal F_{i-1} & & \mathcal F_i
\end{tikzcd}
\tag{4.1.1}\label{eq:a6-02-4-1-1}
\end{equation}
使得 $\operatorname{coker}(f_i^{\leftarrow})$ 在 $\Gamma_{x'_i}$ 上局部自由、秩为 $1$，且 $\operatorname{coker}(f_i^{\rightarrow})$ 在 $\Gamma_{\sigma(x'_i)}$ 上局部自由、秩为 $1$。特别地，$f_i^{\leftarrow}$ 与 $f_i^{\rightarrow}$ 限制到 $X'\times S-\Gamma_{x'_i}-\Gamma_{\sigma(x'_i)}$ 上可逆，且复合
\[
\begin{aligned}
\mathcal F_{i-1}|_{X'\times S-\Gamma_{x'_i}-\Gamma_{\sigma(x'_i)}}
&\xrightarrow{(f_i^{\leftarrow})^{-1}}
\mathcal F^{\flat}_{i-1/2}|_{X'\times S-\Gamma_{x'_i}-\Gamma_{\sigma(x'_i)}}\\
&\xrightarrow{f_i^{\rightarrow}}
\mathcal F_i|_{X'\times S-\Gamma_{x'_i}-\Gamma_{\sigma(x'_i)}}
\end{aligned}
\]
与 $f_i$ 一致。
\end{enumerate}
\end{defn}

对 $X'\times S$ 上的向量丛 $\mathcal F$，记 ${}^\tau\mathcal F:=(\mathrm{Id}_{X'}\times\mathrm{Frob}_S)^*\mathcal F$ 为其 Frobenius 扭。若 $\mathcal F$ 带 Hermite 结构 $h\colon\mathcal F\xrightarrow{\sim}\sigma^*\mathcal F^\vee$，则 ${}^\tau\mathcal F$ 带 Hermite 结构 ${}^\tau h$；当我们谈论「Hermite 丛」${}^\tau\mathcal F$ 时，可以省略这个记号。把 $\mathcal F\in\Bun_{U(n)}(S)$ 视为对象，${}^\tau\mathcal F$ 就是 $\mathcal F$ 在映射 $\mathrm{Frob}\colon\Bun_{U(n)}\to\Bun_{U(n)}$ 下的像。

\ysrc{493}\begin{defn}{4.1.2}设 $r\ge0$ 是整数。我们由笛卡尔图定义 $\mathrm{Sht}^r_{U(n)}$：
\[
\begin{tikzcd}[column sep=large, row sep=large]
\mathrm{Sht}^r_{U(n)} \arrow[r] \arrow[d] & \mathrm{Hk}^r_{U(n)} \arrow[d] \\
\Bun_{U(n)} \arrow[r, "{(\mathrm{Id},\mathrm{Frob})}"'] & \Bun_{U(n)}\times\Bun_{U(n)}
\end{tikzcd}
\]
$\mathrm{Sht}^r_{U(n)}$ 的点称为「秩 $n$ 的 Hermite shtuka」。

具体地，$\mathrm{Sht}^r_{U(n)}$ 的 $S$-点是下列数据的群胚：
\begin{enumerate}
\item $x'_i\in X'(S)$，$i=1,\dots,r$，其图记为 $\Gamma_{x'_i}\subset X'\times S$。它们称为该 shtuka 的 legs。
\item $X'\times S$ 上秩 $n$ 的向量丛序列 $\mathcal F_0,\dots,\mathcal F_r$，各带 Hermite 结构 $h_i\colon\mathcal F_i\xrightarrow{\sim}\sigma^*\mathcal F_i^\vee$。
\item 同构 $f_i\colon\mathcal F_{i-1}|_{X'\times S-\Gamma_{x'_i}-\Gamma_{\sigma(x'_i)}}\xrightarrow{\sim}\mathcal F_i|_{X'\times S-\Gamma_{x'_i}-\Gamma_{\sigma(x'_i)}}$，与 Hermite 结构相容；作为 $X'\times S$ 上底层向量丛的修改，它们与 Hecke 叠 $\mathrm{Hk}^r_{U(n)}$ 定义中的情形一样。
\item 同构 $\varphi\colon\mathcal F_r\xrightarrow{\sim}{}^\tau\mathcal F_0$，与 Hermite 结构相容。
\end{enumerate}
\end{defn}

\begin{rem}{4.1.3}可以证明 $\mathrm{Sht}^r_{U(n)}$ 是局部有限型的 Deligne–Mumford 叠。映射 $\mathrm{Sht}^r_{U(n)}\to(X')^r$ 光滑、分离、等维数，相对维数为 $r(n-1)$。这些以及更多几何性质见 [FYZ24, §6.4]。
\end{rem}

\subsection{特殊闭链}

现在定义 $\mathrm{Sht}^r_{U(n)}$ 上的特殊闭链，它们类似于 Kudla–Rapoport 在酉志村簇上构造的特殊闭链 [KR14]。

\begin{defn}{4.2.1}设 $\mathcal E$ 是 $X'$ 上秩 $m$ 的向量丛。我们定义叠 $Z^r_{\mathcal E}$，其 $S$-点是下列数据的群胚：
\begin{itemize}
\item 秩 $n$ 的 Hermite shtuka $(\{x'_1,\dots,x'_r\},\{\mathcal F_0,\dots,\mathcal F_r\},\{f_1,\dots,f_r\},\varphi)\in\mathrm{Sht}^r_{U(n)}(S)$；
\item $X'\times S$ 上的凝聚层映射 $t_i\colon\mathcal E\boxtimes\OO_S\to\mathcal F_i$，使得同构 $\varphi\colon\mathcal F_r\xrightarrow{\sim}{}^\tau\mathcal F_0$ 把 $t_r$ 与 ${}^\tau t_0$ 缠结，且对每个 $i=1,\dots,r$，修改 $f_i\colon\mathcal F_{i-1}\dashrightarrow\mathcal F_i$ 把 $t_{i-1}$ 与 $t_i$ 缠结，也就是说下图交换：
\[
\begin{tikzcd}[column sep=large, row sep=large]
\mathcal E\boxtimes\OO_S \arrow[r, equal] \arrow[d, "t_0"'] &
\mathcal E\boxtimes\OO_S \arrow[r, equal] \arrow[d, "t_1"'] &
\cdots \arrow[r, equal] \arrow[d] &
\mathcal E\boxtimes\OO_S \arrow[r, "\sim"] \arrow[d, "t_r"'] &
{}^\tau(\mathcal E\boxtimes\OO_S) \arrow[d, "{}^\tau t_0"'] \\
\mathcal F_0 \arrow[r, dashed, "f_1"] &
\mathcal F_1 \arrow[r, dashed, "f_2"] &
\cdots \arrow[r, dashed, "f_r"] &
\mathcal F_r \arrow[r, "\sim"] &
{}^\tau\mathcal F_0
\end{tikzcd}
\]
\end{itemize}
在下文中，写这类图时我们通常省略记号中的「$\boxtimes\OO_S$」因子。
\end{defn}

\begin{rem}{4.2.2}存在显然的映射 $Z^r_{\mathcal E}\to\mathrm{Sht}^r_{U(n)}$，投影到前一项所列的数据。由 [FYZ24, Proposition 7.5]，该映射有限，故诱导 Chow 群上的前推映射。
\end{rem}

\ysrc{494}\begin{defn}{4.2.3}设 $A_{\mathcal E}(\F_q)$ 是满足 $\sigma^*a^\vee=a$ 的 Hermite 映射 $a\colon\mathcal E\to\sigma^*\mathcal E^\vee$ 构成的 $\F_q$-向量空间。

设 $(\{x'_i\},\{\mathcal F_i\},\{f_i\},\varphi,\{t_i\})\in Z^r_{\mathcal E}(S)$。由 $Z^r_{\mathcal E}$ 的定义中 $t_i$ 之间的相容性，序列
\begin{equation}
\mathcal E\boxtimes\OO_S\xrightarrow{t_i}\mathcal F_i\xrightarrow{h_i}\sigma^*\mathcal F_i^\vee\xrightarrow{\sigma^*t_i^\vee}\sigma^*\mathcal E^\vee\boxtimes\OO_S
\tag{4.2.1}\label{eq:a6-02-4-2-1}
\end{equation}
中映射的复合对每个 $i$ 都一致，并且 $i=r$ 时的 (4.2.1) 也与 $i=0$ 时的 (4.2.1) 的 Frobenius 扭一致。因此复合映射 (4.2.1) 对每个 $i$ 都给出 $A_{\mathcal E}(S)$ 中同一个元素，而且它被 Frobenius 固定，故必来自 $A_{\mathcal E}(\F_q)$。这定义了映射 $Z^r_{\mathcal E}\to A_{\mathcal E}(\F_q)$。对 $a\in A_{\mathcal E}(\F_q)$，记 $Z^r_{\mathcal E}(a)$ 为 $Z^r_{\mathcal E}$ 在 $a$ 上的纤维。有
\[
Z^r_{\mathcal E}=\coprod_{a\in A_{\mathcal E}(\F_q)}Z^r_{\mathcal E}(a).
\]
$Z^r_{\mathcal E}$ 或 $Z^r_{\mathcal E}(a)$ 称为余秩 $m$（带 $r$ 个 legs）的特殊闭链。
\end{defn}

\subsection{特殊闭链的虚基本类}

受 Kudla 纲领启发，我们期待对每个特殊闭链 $Z^r_{\mathcal E}(a)$（$a\in A_{\mathcal E}(\F_q)$）配置 Chow 群 $\mathrm{CH}_*(\mathrm{Sht}^r_{U(n)})$ 中的一个类 $[Z^r_{\mathcal E}(a)]$，使得以 $[Z^r_{\mathcal E}(a)]$ 为 Fourier 系数的 Fourier 级数是「自守的」。我们把这样的自守形式称为高阶 theta 函数，因为当 $r=0$ 时它回到经典 theta 函数的定义。

然而稍加审视即可发现：当 $r>0$ 时，闭链 $Z^r_{\mathcal E}(a)$ 的维数必随 $a$ 变化。例如，若 $a=0$，则 $Z^r_{\mathcal E}(0)$ 显然满射到 $\mathrm{Sht}^r_{U(n)}$ 上。另一方面，不难看出当 $r>0$ 且 $a\ne0$ 时，$Z^r_{\mathcal E}(a)$ 的维数比 $\mathrm{Sht}^r_{U(n)}$ 小。因此，在高阶 theta 函数的定义中不应使用朴素的闭链类 $[Z^r_{\mathcal E}(a)]^{\mathrm{naive}}$；我们必须构造虚基本闭链 $[Z^r_{\mathcal E}(a)]^{\mathrm{vir}}$，它们都应当落在 $\mathrm{CH}_{(n-m)r}(\mathrm{Sht}^r_{U(n)})$ 中，其中 $m=\operatorname{rank}\mathcal E$。

\subsubsection{非奇异项}

有许多启发式指导这样的构造。一个关键输入是：当 $m=\operatorname{rank}\mathcal E=1$，且 $a$ 非奇异——即 $a\colon\mathcal E\to\sigma^*\mathcal E^\vee$ 作为凝聚层映射是单射（等价地，在 $X'$ 的一般点上是同构）——时，$Z^r_{\mathcal E}(a)$ 是「正确」余维（即 $r$）的 LCI。此时，启发式（事后由 §3.2.1 证实）提示可以取 $[Z^r_{\mathcal E}(a)]^{\mathrm{vir}}=[Z^r_{\mathcal E}(a)]^{\mathrm{naive}}$。

更一般地，利用这一观察，可以按 [KR14] 的原始方法对任意非奇异 $a$ 定义 $[Z^r_{\mathcal E}(a)]^{\mathrm{vir}}$（要点在于允许 $m>1$）：把 $Z^r_{\mathcal E}(a)$ 表示为「$m=1$ 且系数标号非奇异」的特殊闭链的「导出交」的一个开闭分支。这一情形下 $[Z^r_{\mathcal E}(a)]^{\mathrm{vir}}$ 的定义在 [FYZ24, §7] 中给出。

\subsubsection{奇异项}

奇异 $a$ 的 $[Z^r_{\mathcal E}(a)]^{\mathrm{vir}}$ 的定义需要更多启发。在数域上，类比问题的答案相对简单：此时 $a$ 可以表示成形如 $a_{\mathrm{ns}}\oplus0$ 的矩阵，其中 $a_{\mathrm{ns}}$ 是非奇异的、秩 $m'\le m$。于是相应的虚基本类等于 $a_{\mathrm{ns}}$ 的虚基本类乘以某个「tautological」线丛的 Chern 类的 $(m-m')$ 次幂。

但在手头的函数域情形，答案要复杂得多：它一般是无穷多项之和，每一项涉及不同的「tautological」丛的最高 Chern 类。我们已经可以说明为什么在特殊的

% 本句在印刷 494 页末中断于 “in the special”，其后半句续于下一 chunk 开头。
